% Source: https://artofproblemsolving.com/community/c3h3633198_usa_tst_for_vimc_2025_a3_quirky_nt
% Stable identity: aops-topic-3633198-q01
Every integer below $50$ has at most three distinct prime factors. If $p^a\mid n$ exactly with $a\ge 2$, the fraction of divisors divisible by $p$ is $\dfrac{a}{a+1}$. These cannot be adjacent, so their number is at most one more than the number not divisible by $p$: $(a-1)\tau(\dfrac{n}{p^a})\le 1$. Thus either $n=p^2$, or no prime exponent exceeds one. If a squarefree $n$ has three primes, it is $30$ or $42$. In a coprime-divisor path, $n$ is adjacent only to $1$, so is an endpoint. Each of $\dfrac{n}{p},\dfrac{n}{q},\dfrac{n}{r}$ is adjacent only to $1$ and its missing prime. At most one can use the remaining edge from $1$; the other two must be endpoints, impossible since only one endpoint remains. Finally $1$, any prime $p$, a prime square $p^2$, and a product $pq$ of distinct primes work, with orders $(1)$, $(1,p)$, $(p,1,p^2)$, and $(pq,1,p,q)$. Primes below $50$ sum to $328$; prime squares are $4,9,25,49$, summing to $87$; distinct-prime products are $6,10,14,15,21,22,26,33,34,35,38,39,46$, summing to $339$. The answer is $1+328+87+339=\boxed{755}$.
